% ED26-BM-RAT2021-8-5: complete Hertwig-Kozyreva 2021 source treatment.
\chapter{Bounded Rationality: Principles, Procedures and Environments}
\label{sec:p17-rationality-8-5}

Bounded rationality examines how agents achieve their goals with limited
time, knowledge, foresight and computational resources in environments
that may exceed their capacity to describe or search. Hertwig and
Kozyreva distinguish the principle that these bounds matter from a
positive theory of how choices are actually produced. They reconstruct
Simon's approach and compare optimization under constraints, heuristics
and biases, and ecological rationality
\cite[chapter~8.5, pp.~505--515]{HertwigKozyreva2021Bounded}.
These programs give different roles to optimization, cognitive mechanisms
and environmental structure. Their common concern with limits does not
make their explanations or standards identical.

\section{From a principle of limits to a theory of choice}
\label{sec:p17-br-foundations}

\begin{hbdefinition}{def:p17-br-principle-theory}{Principle, positive theory and prescription}
As a principle, bounded rationality directs attention to the mismatch
between the agent's capacities and the problem's demands. It identifies
time, information, foresight and computation as relevant bounds.
As a positive theory, it must specify and test mechanisms of search,
representation and selection that can operate within those bounds.
A prescriptive use additionally asks which available procedures help the
agent achieve a stated goal in a stated environment. A descriptive fit
alone does not establish that a procedure ought to be used.
\end{hbdefinition}

In Simon's early organizational work, limits could be described as
departures from omniscient choice. The chapter locates a subsequent
change in his 1955 and 1956 work: the task becomes explaining behavior
under the conditions organisms actually encounter, then developing a
different account of rational choice. This is more than deleting details
from an ideal maximizing agent. The psychological process and the
environment in which it acts become objects of study
\cite[pp.~505--506 and note~1]{HertwigKozyreva2021Bounded}.

\begin{hbdefinition}{def:p17-br-three-rationalities}{Substantive, procedural and ecological rationality}
Substantive rationality evaluates the choice against a specified goal,
constraints and outcome criterion. In the classical models considered by
the chapter, its standards include consistent preferences and expected
utility. Procedural rationality evaluates and explains the processes by
which an agent arrives at a choice, including the information and
computation actually used. Ecological rationality evaluates a procedure
in relation to the structure of its environment and its success at the
agent's goals. The latter asks when a procedure's architecture fits the
information distribution and demands of a task.
\end{hbdefinition}

The chapter collects Simon's terms for his classical target: global,
perfect, objective, substantive and full rationality, the rational
\emph{homo economicus}, and the Olympian model. Their historical target
includes expected and subjective expected utility with coherence standards
such as completeness and transitivity. They are not all definitions of
one modern mathematical object. Ecological rationality here is the
psychological research program described by these authors; the chapter's
note~6 expressly distinguishes Vernon Smith's different economic use of
the same phrase \cite[pp.~506, 513]{HertwigKozyreva2021Bounded}.

Simon's scissors image joins two parts of an explanation: the task
environment and the actor's computational capacities. His labyrinth
image describes decisions as successive choices on a branching path,
with future possibilities discovered along the way. Neither image by
itself is an algorithm or a proof that a particular heuristic succeeds.
A model must specify what is encountered, what information is available,
which operations are possible, and what terminates the process.

\section{The assumptions and limits of the Olympian model}
\label{sec:p17-br-olympian}

The chapter's criticism addresses the demands of a familiar complete
choice description: mutually exclusive and exhaustive relevant states,
specified consequences, comparable values for those consequences, and
known outcomes or definite probabilities. The agent must then identify a
maximizer over the alternatives. Such a description can be useful in a
closed task. It leaves little room for consequences that were not
anticipated, alternatives not yet encountered, or goals that cannot be
put into a complete ordering. The authors connect this contrast with
Savage's distinction between simplified small worlds and the uncertain
conditions of ordinary choice \cite[pp.~506--507]{HertwigKozyreva2021Bounded}.

The finite expected-utility theorem in \cref{thm:p17-finite-vnm} makes
the preference assumptions explicit. It does not show that an observed
person supplies those preferences, knows all the relevant consequences,
or has an executable procedure for discovering them. Similarly,
\cref{thm:p17-sequential-backward} proves an optimal policy for the
finite controlled model in \cref{def:p17-sequential-model}; it starts
with its states, actions, transition probabilities and rewards. An
existence or representation theorem and the process of supplying and
using its inputs answer different questions.

The chapter's firefighter example concerns expert action under time
pressure, serious consequences and uncertainty. The expert recognizes a
situation from experience, generates a plausible action, mentally checks
it, and modifies or replaces it if that check fails. The case illustrates
recognition and sequential testing of an action, rather than generation
and comparison of every possible action. It is attributed to the expert
decision account discussed through Klein; it is not a newly measured
frequency or an assertion that every expert follows one invariant rule.
The chapter also names emergency-room doctors, police officers, soldiers
and chess players as contexts in which limits matter. Transfer between
these occupations requires its own task and evidence record.

The accompanying uncertainty distinction matters. Epistemic uncertainty
concerns missing knowledge; aleatory uncertainty concerns stochastic
features of the environment. The chapter relates it to the epistemic
and frequency interpretations historically associated with probability.
Better information can change the first without eliminating the second.
An exhaustive outcome model may itself be unavailable, so uncertainty
is not always represented adequately by attaching a known probability
to each member of an already complete list
\cite[pp.~506--507 and note~3]{HertwigKozyreva2021Bounded}.

Chess supplies a separate illustration of search growth. Under the
chapter's schematic assumptions of thirty options at each of eighty
successive individual moves, there are \(30^{80}\), approximately
\(1.48\times10^{118}\), paths. This is a branching illustration,
not an exact count of legal positions: branching varies, paths can
terminate, and different paths can reach a common position. The practical
success of selective search does not require exhaustive enumeration of
that schematic tree. The 1997 Deep Blue system used a parallel search
architecture, evaluation features, search extensions and a game database.
Its engineers report an average of 126 million positions per second in
match searches exceeding one minute, with a maximum sustained speed of
330 million; tactical and quiet positions had different typical speeds
\cite[pp.~57--60]{CampbellHoaneHsu2002DeepBlue}.
Throughput and selective search should therefore be specified separately.
Neither a peak speed nor the game-tree illustration is a general
computational-impossibility theorem.

The chapter then broadens the comparison: social situations can have
unclear rules, many participants, numerous possible actions and conflicting
goals. Formalizing them as chess-like closed games can omit precisely
the information discovery or uncertainty that governs the actual choice.
This does not invalidate a theorem whose supplied model is correct; it
asks whether that model and its procedure describe the situation at hand.
The historical discussion of Simon's 1979 Nobel lecture emphasizes
uncertainty and imperfect competition as settings in which he argued
that the classical account was inadequate. This is his criticism of
model adequacy as reconstructed by the chapter, rather than a newly
proved impossibility claim \cite[note~4, p.~513]{HertwigKozyreva2021Bounded}.

\section{Four principles of Simon's investigation}
\label{sec:p17-br-four-principles}

The first principle requires a behaviorally informed theory of the process.
An explanation examines both the final alternative and how it was found.
The chapter contrasts Simon's insistence on cognitive process realism
with Friedman's emphasis on whether a theory's predictions suffice for
their purpose. Prediction and process explanation can support different
model comparisons; an adequate endpoint prediction need not identify the
mechanism that produced it \cite[pp.~507--508]{HertwigKozyreva2021Bounded}.

Second, the methods may be approximate. That need not mean that their
solutions are inferior or their cognition unsophisticated. Recognition,
emotions, motion tracking and adaptive forgetting can exploit substantial
learned or evolved capacities while requiring little computation at the
moment of choice. The process must still specify how those capacities
apply and under which environmental conditions they help.

Third, satisficing uses heuristic search with adjustable aspirations.
\begin{hbdefinition}{def:p17-br-satisficing}{Search and an aspiration rule}
A satisficing procedure specifies how alternatives are encountered, an
aspiration that counts as acceptable, a rule for adjusting that aspiration,
and a stopping rule. It searches for a satisfactory alternative rather
than requiring proof that the chosen alternative is best in an exhaustively
known set. With a time or inspection limit, it also requires a feasible
fallback if the aspiration is never met. Whether earlier alternatives
remain available is part of the model.
\end{hbdefinition}

This is a procedure for finding options, not merely a relabelling of an
already known second-best option. Aspiration can change after unsuccessful
search. The chapter's example of sequentially encountered potential
partners emphasizes that an unchosen alternative may disappear. The
order and availability of alternatives then matter. Satisficing is one
bounded procedure among others, not a necessary feature of every
boundedly rational act.

Fourth, explanation includes the environment. Simon's ant-on-a-beach
example, shown in the source's Figure~8.5.1, contrasts a complicated
observed route with a simple organism interacting with stones, puddles
and other obstacles. A local avoidance response can generate a circuitous
path. Inferring an elaborate cognitive mechanism from the path alone
confounds process with surroundings. The chapter explicitly allows that
real ants have other determinants of action; the illustration is a
mechanism argument, not a quantified ant experiment.

\section{Optimization under constraints}
\label{sec:p17-br-optimization}

The first research program retains optimization while incorporating limits.
The chapter describes Stigler's treatment of costly search: continue
looking for alternatives or information while the expected marginal
benefit exceeds the marginal cost. A car buyer may stop when the expected
gain from a better option no longer justifies time and effort. Its
historical note~5 identifies Good's earlier discussion of information
and decision costs in 1952. The program is consequently not represented
as an innovation of Stigler alone \cite[pp.~508--509, 513]{HertwigKozyreva2021Bounded}.

This approach can be executable in a closed model with known search
returns and costs. The objection concerns the work needed to obtain
those quantities: evaluating the opportunities forgone, the distribution
of undiscovered options, or the benefit of more computation can itself
be difficult. Adding these quantities to an objective does not supply
their values or make their optimization computationally feasible.
Simon's criticism is that putting search back inside utility maximization
can obscure a principal problem his theory addresses: choice with
incomplete information without executing an optimization procedure.

Accordingly, the program's criterion and its implementation should be
reported separately. A resource-constrained optimum is one candidate
benchmark. A psychologically realistic model still requires the actual
available information, processing steps and stopping behavior. Costs
are not necessarily known, and their mere inclusion is not evidence that
the agent has overcome a resource bound.

\section{Heuristics and biases}
\label{sec:p17-br-heuristics-biases}

The second program compares judgments and decisions with standards from
probability, statistics and rational choice. Its influential map of
departures includes base-rate errors, overconfidence, conjunction
judgments, hot-hand judgments, and failure to account for regression
or aggregation. Representativeness, availability, and anchoring with
adjustment are mechanisms proposed to explain families of findings.
The original program does not say that these shortcuts always fail:
they often work, while particular task conditions can produce systematic
errors \cite[pp.~509--510]{HertwigKozyreva2021Bounded}.

The chapter describes attribute substitution through three conditions:
the target property is comparatively inaccessible, a substitute property
is readily available, and a controlling process does not reject its use.
Representativeness can substitute resemblance for a probability judgment.
In the Linda illustration, resemblance to the supplied profile may favor
the description of a feminist bank teller over that of a bank teller.
For two events interpreted on the same probability space, however,
the conjunction cannot have probability exceeding that of the bank-teller
event. Interpreting and measuring a response is a separate question;
the actual task and language distinctions are developed in
\cref{sec:p19-conjunction-judgments,sec:p19-conjunction-language}.

The chapter's account of Kahneman additionally uses two fictional systems
as explanatory devices. System~1 supplies fast, associative responses;
System~2 can reflect on and correct them but is effortful and does not
always do so; the account includes fatigue as a limit on its operation.
These model descriptions do not identify two independently
verified anatomical modules or prove that a given response came from one
of them. Evidence must distinguish mechanism predictions from a convenient
description of the result.

Prospect theory is presented as an influential descriptive revision of
expected-utility choice, using reference dependence, probability weighting
and loss aversion. In its original risky-gamble setting, the possibilities
and probabilities are supplied. Retaining expected utility as a normative
benchmark while changing its descriptive predictions differs from
Simon's proposal to explain choice when the complete risky-choice
description is unavailable. A theory for supplied lotteries does not
automatically provide a procedure for undiscovered consequences.

The chapter gives three lines of objection to this program. Content-free
coherence norms can offer a narrow view of cognitive success. Responses
can change when information is presented in different formats, so an
apparent illusion need not be impervious to instruction or representation.
Finally, a procedure adapted to another environment can be misdiagnosed
in an experiment whose structure differs from that environment. These
are objections to interpretation and transfer, not exceptions to the
probability inclusion theorem. Claims about format effects or adaptation
need actual task evidence rather than inference from an error label.

\section{Ecological rationality}
\label{sec:p17-br-ecological}

The third program treats a heuristic as a strategy that leaves some
information unused in pursuit of speed, frugality, accuracy or another
stated goal. It investigates an adaptive collection of tools learned
individually, culturally or through evolution. Its proposed correspondence
standards evaluate success in the world; the relevant criterion may include
accuracy, speed, robustness or accountability. Choosing such a criterion
is a normative step, while measuring performance under it is an empirical
or mathematical task \cite[pp.~510--512]{HertwigKozyreva2021Bounded}.

Competitive evaluation is central. Specify alternative strategies and a
reference class of environments, rather than judging a heuristic in isolation.
A strategy may fit one environment and fail badly in another. The chapter
draws three consequences: simplicity does not necessarily incur inferior
performance; environmental structures and the processes producing them
require their own theory; and successful heuristics can operate across
judgment, inference, classification and choice under uncertainty. Its
risk--reward example points to foraging theory as a possible source of
structured environmental explanations, rather than merely compiling
unrelated task descriptions.

The conventional accuracy--effort story assumes that less processing saves
effort at some loss of accuracy. Ecological work investigates cases in which
less information or computation instead preserves or improves prediction.
This is a conditional performance claim. It does not imply that ignoring
information is always beneficial or that any simple rule outperforms an
appropriately specified complex rule.

Fitting and prediction must be distinguished. For nested model classes and
the same globally minimized training loss, adding flexibility cannot worsen
the best achievable training fit. Held-out prediction has no corresponding
monotonic guarantee. Parameter count alone does not establish nestedness,
successful optimization, or future performance. The chapter's account
invokes an often observed tradeoff between too little and too much flexibility
when estimated parameters depend on small or unreliable samples.

\begin{hbtheorem}{thm:p17-br-bias-variance}{proved=annex:p17.br.bias-variance}{A specified squared-error mechanism for a less-is-more effect}
At a fixed prediction input, let \(Y=m+\epsilon\), where
\(E\epsilon=0\), \(E\epsilon^2=\sigma^2<\infty\). Let a trained
prediction \(\widehat m\) have finite second moment and be independent
of the test noise. Expectations here use the stipulated joint training
and test law. Then
\[
 E[(\widehat m-Y)^2]=(E\widehat m-m)^2+
                         \operatorname{Var}(\widehat m)+\sigma^2.
\]
A fixed prediction \(a\) has lower risk than an unbiased trained one
with variance \(v\) exactly when \((a-m)^2<v\).
For stipulated \(a=0\), \(\sigma^2=1\), \(v=4\), the fixed
prediction has risk one against five when \(m=0\), but risk ten
against five when \(m=3\). These are analytical examples, not sample
estimates or claims that complexity uniquely determines bias or variance.
\end{hbtheorem}

Another mechanism is noncompensation in cue integration. A correlation
between a cue and the criterion by itself is not a condition that the cue
outweighs all others. The precise binary linear rule used by Şimşek has
noncompensatory weights when each weight exceeds the sum of all later
weights in the specified ordering
\cite[§2.3]{Simsek2013LinearDecision}.
For a negative coefficient on a binary cue, replacing the cue by its
complement gives a positive coefficient and changes only the intercept;
that intercept cancels in a paired comparison. Zero-weight cues can be
discarded. Thus cue direction and ordering remain explicit inputs.

\begin{hblemma}{thm:p17-br-noncompensation}{proved=annex:p17.br.noncompensation}{When a first discriminating cue preserves a linear comparison}
For binary attributes and nonnegative weights \(w_1,\ldots,w_k\),
where \(k\geq1\) is an integer, suppose
\(w_i>\sum_{j>i}w_j\) for every \(i<k\).
Let \(d_i=x_i(a)-x_i(b)\in\{-1,0,1\}\).
If a positive-weight cue discriminates, the sign of \(\sum_iw_id_i\)
equals the sign of the first such \(d_i\). If none discriminates,
the linear criterion ties. A matching lexicographic rule, omitting
zero-weight cues, therefore makes the same
comparison; its cue order and directions are assumptions of the result.
\end{hblemma}

The rule is a mathematical performance guarantee under those conditions,
not a theorem about every instance of Take-the-Best or LEX. Those
strategies additionally require a way to order cues and learn their
directions. A cue order imported from another environment can fail.

As a bounded empirical grounding, Şimşek's study evaluates computational
paired comparisons in 51 selected datasets with 12--601 objects and
3--21 attributes. These are datasets, not 51 sampled human populations.
The analysis uses numeric and median-split binary versions, scaled cues,
mean replacement of missing attributes and exclusion of missing criterion
values. Its means, median splits and ranges use the dataset's objects;
they are not presented as preprocessing estimated wholly separately from
the compared objects. Elastic-net coefficients are selected by training-set cross
validation; paired comparisons use leave-pair-out training for smaller
datasets or a 1,000-fold pair allocation for larger ones, with five
random-seed replications. Its noncompensation measure asks whether the
first discriminating cue, in the fitted linear rule's order, agrees with
that rule's full comparison. The reported across-dataset medians are
0.85 for numeric versions and 0.96 for binary versions
\cite[§4, pp.~4--7]{Simsek2013LinearDecision}.
This is rule agreement under the specified evaluation, not accuracy at
predicting the target, not the fraction of people using a heuristic, and
not a guarantee for another cue order. Paired comparisons share objects,
and seed replications do not create independent populations. The paper
does not supply a population sampling confidence interval for those
medians. Its discussion expressly separates the possibility that matching
heuristics succeed from a finding that an actual psychological heuristic
necessarily uses that matching cue order.

The program also develops decision aids, including fast-and-frugal trees.
The chapter names geographic profiling, clinical prescribing and marketing
allocation as application areas and links this work to boosting competences.
It contrasts this emphasis on usable tools with nudging based on predictable
choice tendencies. Neither an application's name nor a theoretical
comparison supplies its deployment performance. Each aid needs a defined
task, feasible actions, observations, criterion and evidence of transfer.
The authors characterize the heuristics-and-biases program chiefly as
mapping departures and the ecological program as additionally constructing
new usable heuristics. This comparison of research priorities does not
assert that representation changes or other remedial work are absent
from the former program.

\section{A feasible-choice and environment contrast}
\label{sec:p17-br-worked-environments}

\begin{hbexample}{def:p17-br-cue-instance}{Two alternatives under a one-inspection deadline}
An agent chooses between two named alternatives. For the analyst their
attribute vectors are \(a=(1,0,0)\) and \(b=(0,1,1)\); the agent
initially has the labels without those revealed coordinate values.
Inspecting a cue reveals that coordinate of both alternatives and costs
one time unit. The deadline permits one inspection. The fixed strategy
inspects cue~1 and chooses its higher-valued alternative; a comparison
that explicitly acquires and sums all three coordinates takes three
time units and is infeasible under this deadline.

For evaluation, specify two environments: \(N\) scores an alternative
by weights \((4,2,1)\), while \(C\) uses \((1,2,4)\).
These are stipulated outcome criteria known to the analyst. Their full
values need not be available to the acting agent. The example does not
assume that each environment is a fixed stochastic population or a
sample of human decisions.
\end{hbexample}

\begin{hbtheorem}{thm:p17-br-environment-contrast}{proved=annex:p17.br.environment-contrast}{The same feasible procedure can succeed or fail}
In environment \(N\), the cue-1 strategy chooses \(a\), whose
score is four versus three for \(b\). In environment \(C\), it
still chooses \(a\), whose score is one versus six for \(b\).
The same one-inspection procedure thus matches the criterion in the first
environment and fails in the second. A cue-3 strategy is also feasible
and succeeds in \(C\); obtaining the new cue ordering is an additional
learning or design task. The explicit full-sum comparator's infeasibility
does not imply that no other feasible procedure can choose correctly.
\end{hbtheorem}

This comparison assesses correspondence with a specified environment.
It does not show a failure of transitivity or completeness in
\cref{def:p17-weak-preference}. It also differs from an already solved
expected-utility maximization problem: the observation and processing
procedure is part of the feasible-choice account. A representation of
the desired ordering does not make its discovery free.

For a satisficing contrast, suppose offers arrive in sequence and rejected
offers cannot be recalled. Inspect at most two offers, each at unit cost,
accept the first with value at least seven, and otherwise accept the final
offer. For the stipulated sequence \((6,8)\), the procedure takes eight
after two inspections, giving net value six. A threshold of five instead
accepts six after one inspection, giving net value five. For sequence
\((6,4)\), the threshold-seven procedure gives net value two, while
the threshold-five procedure still gives five. These direct calculations
make aspiration, search order, disappearance and fallback explicit.
They do not assert that either threshold is optimal under an unspecified
distribution of future offers.

\section{Evidence, disagreements and continuing questions}
\label{sec:p17-br-evidence-limits}

An empirical evaluation should record the processes as well as outcomes
when its claim concerns procedural rationality. Use
\cref{def:p19-evidence-record} to specify population or dataset selection,
task, available aids, timing, response measure, uncertainty and comparison
criterion. The chapter's firefighter narrative, ant illustration, Linda
example, model comparisons and policy applications have different evidence
roles. A stipulated environment proves a conditional possibility; it does
not measure how often humans occupy that environment or use its rule.
Predictive success, a credible mechanism and a normative recommendation
need their own arguments.

The chapter closes with unresolved research directions as understood in
2021: a theory-based taxonomy of environmental structures; the ways agents'
heuristics jointly produce environmental complexity; a broader collection
of cognitive tools beyond simple heuristics; and integration of search,
learning and heuristic choice. Its parental-investment and micromotive
references motivate the possibility that simple individual rules generate
complex aggregate patterns. These are research questions and conceptual
examples, not assertions that a unified framework has already been
established or that the same questions have unchanged present status
\cite[p.~512]{HertwigKozyreva2021Bounded}.

For consultation, specify the goal and benchmark, the resources actually
available, the search and stopping process, and the environmental conditions
on which its success depends. State whether the claim describes observed
behavior, predicts it, proves a formal guarantee, or recommends a feasible
procedure. Keeping those claims explicit preserves the source's central
point: understanding choice requires both the actor's capacities and the
world in which the actor uses them.
